\section{Results}
\label{sec:results}

We report 45 configurations: three probe-side thicknesses
($P \in \{\PThin, \PMid, \PThick\}$ rows on the hot key) crossed with six
strategies, a salt sweep $k \in \{1,\dots,32\}$, and --- for the AQE arms --- a
paired run with partition coalescing disabled. Each cell runs one discarded
warmup and five timed repetitions, in randomised order. Every number below is
generated from the results file by the released analysis code; none is
transcribed by hand.

\subsection{The measurement floor}

Selective salting at $k=1$ compiles to exactly the same physical plan as the
unmitigated baseline, so the difference between them is pure measurement noise.
We report it first, because it bounds what any other number in this section can
mean:

\begin{center}
\begin{tabular}{lrrr}
\toprule
$P$ & \PThin & \PMid & \PThick \\
\midrule
Baseline wall (s)                & \BaseWallThin & \BaseWallMid & \BaseWallThick \\
Identical-plan control (s)       & 3.09 & 6.97 & 19.84 \\
\textbf{Measurement floor}       & \textbf{18.3\%} & \textbf{0.1\%} & \textbf{6.5\%} \\
\bottomrule
\end{tabular}
\end{center}

The floor varies sharply with run length, and this governs how the rest of the
section may be read. At $P=\PMid$ the two identical plans agree to 0.1\%, and a
difference of a few percent is meaningful. At $P=\PThin$ they differ by 18.3\%:
those runs last under three seconds, short enough that JVM and scheduler noise
dominates, and \textbf{no claim can be supported at that probe thickness at
all}. We report the $P=\PThin$ column throughout for completeness and draw no
conclusions from it.

Two points about how this number is produced. All 45 cells ran in a single
Spark application: an earlier attempt gap-filled four cells from a second
application, and those cells were systematically slower and drifting within
their own run, which inflated this control from 2\% to 29\%. Wall-clock
comparisons across Spark applications measure JVM warmth, not plans, and the
harness now records the application id and refuses to compare across it.
Separately, an earlier version executed cells in configuration order, placing
the baseline first on a cold JVM; that inflated the same control to 17.7\% and,
with it, every speedup measured against that baseline. Cell order is now
randomised under a fixed seed.

\subsection{The baseline is severely skewed}

Without mitigation, the join stage's longest task takes \BaseMaxTaskThin~s,
\BaseMaxTaskMid~s and \BaseMaxTaskThick~s at the three probe thicknesses, against
median tasks of a fraction of a second. At $P=\PThick$ a single task runs
roughly two orders of magnitude longer than its median peer. This is a textbook
straggler, and no amount of additional compute would remove it.

\subsection{AQE did not help at any probe thickness}

This is the paper's central measurement, and it is a null result.

\begin{center}
\begin{tabular}{lrrr}
\toprule
$P$ & \PThin & \PMid & \PThick \\
\midrule
Baseline wall (s)            & \BaseWallThin & \BaseWallMid & \BaseWallThick \\
AQE wall (s)                 & \AqeWallThin & \AqeWallMid & \AqeWallThick \\
\textbf{AQE speedup}         & \textbf{\AqeSpeedupThin$\times$} & \textbf{\AqeSpeedupMid$\times$} & \textbf{\AqeSpeedupThick$\times$} \\
\midrule
Baseline longest task (s)    & \BaseMaxTaskThin & \BaseMaxTaskMid & \BaseMaxTaskThick \\
AQE longest task (s)         & \AqeCoalMaxTaskThin & \AqeCoalMaxTaskMid & \AqeCoalMaxTaskThick \\
\bottomrule
\end{tabular}
\end{center}

At $P=\PMid$, where the measurement floor is 0.1\%, AQE returns
\AqeSpeedupMid$\times$. At $P=\PThick$, floor 6.5\%, it returns
\AqeSpeedupThick$\times$ --- inside the floor. And the row that settles it is
the last one: enabling Adaptive Query Execution on a join whose longest task
ran \BaseMaxTaskThick~s produced a longest task of \AqeCoalMaxTaskThick~s.
\textbf{The critical path is not shortened at any probe thickness; in all three
it is marginally longer.}

We are careful about what this does and does not say. It is not a claim that
AQE is useless in general --- it does other useful things, and one of them
shows up below. It is a claim that on \emph{this} join, with Spark's skew
machinery enabled and a straggler two orders of magnitude out, AQE delivered no
improvement distinguishable from measurement noise, and left the straggler
itself untouched. The next two subsections establish why, and the answer is
specific enough to act on.

\subsection{The benefit that does exist is coalescing, not skew handling}

AQE bundles several rewrites. To separate them we ran every AQE cell twice,
identical but for \texttt{spark.sql.adaptive.coalescePartitions.enabled}:

\begin{center}
\begin{tabular}{lrrr}
\toprule
$P$ & \PThin & \PMid & \PThick \\
\midrule
AQE, coalescing on (s)   & \AqeCoalWallThin & \AqeCoalWallMid & \AqeCoalWallThick \\
AQE, coalescing off (s)  & \AqeNoCoalWallThin & \AqeNoCoalWallMid & \AqeNoCoalWallThick \\
Baseline (s)             & \BaseCoalWallPreciseThin & \BaseCoalWallPreciseMid & \BaseCoalWallPreciseThick \\
\midrule
Longest task, coal.\ off (s) & \AqeNoCoalMaxTaskThin & \AqeNoCoalMaxTaskMid & \AqeNoCoalMaxTaskThick \\
\bottomrule
\end{tabular}
\end{center}

With coalescing disabled, AQE is indistinguishable from no AQE at all. The
longest task is unchanged in every column. Whatever small benefit AQE provided
came from merging small post-shuffle partitions --- which lowers fixed
per-task overhead and does precisely nothing for a straggler.

\subsection{The skew-join rule never fired, and why}

Reading the \emph{final} adaptive plan for every AQE cell explains the result
directly. Capturing it requires care: executing into a \texttt{noop} sink builds
a separate query execution and leaves the inspected plan reporting
\texttt{isFinalPlan=false}, showing the pre-adaptive shape and hiding exactly
the rewrite one wants to audit. Read after an action on the DataFrame itself,
every plan contains \texttt{AQEShuffleRead coalesced} and never
\texttt{AQEShuffleRead skewed}. \textbf{Spark's skew-join rule did not fire.}

The reason is a precondition that is documented but almost never discussed. A
skewed partition is split into pieces of roughly
\texttt{advisoryPartitionSizeInBytes}; if the hot partition is \emph{smaller}
than the advisory size, there is no split to make and the rule is a no-op
however skewed the partition is. The measured hot shuffle partition here is
\MeasHotMB~MB against a \NPART-partition advisory size of 16~MB.

This is a claim about a non-event, so it needs a positive control. We swept the
two documented trigger conditions to their most permissive settings ---
\texttt{skewedPartitionThresholdInBytes} from 64~MB down to 256~KB, and
\texttt{skewedPartitionFactor} from 5 down to 1. The rule still never fired.
Lowering only the advisory size engaged it immediately:

\begin{center}
\begin{tabular}{lrrl}
\toprule
Partitions & Advisory & Threshold & Skew-join applied \\
\midrule
16  & 16\,MB   & 8\,MB    & no  \\
16  & 1\,MB    & 1\,MB    & \textbf{yes} \\
64  & 1\,MB    & 1\,MB    & \textbf{yes} \\
200 & 256\,KB  & 256\,KB  & \textbf{yes} \\
\bottomrule
\end{tabular}
\end{center}

One configuration change, no code change. We regard this as the most
practically consequential finding in the paper: practitioners tune the two
parameters whose names contain the word \emph{skew}, and the parameter that
actually gates the optimisation at moderate data volumes is the one that does
not. Section~\ref{sec:decision}'s procedure encodes all three preconditions for
this reason, and predicted the non-event from data statistics alone.

\subsection{Salting worked, and its advantage grows with the probe side}

\begin{center}
\begin{tabular}{lrrr}
\toprule
$P$ & \PThin & \PMid & \PThick \\
\midrule
Best selective-salt $k$        & \SaltBestKThin & \SaltBestKMid & \SaltBestKThick \\
Salting speedup                & \SaltSpeedupThin$\times$ & \SaltSpeedupMid$\times$ & \SaltSpeedupThick$\times$ \\
Longest task, baseline (s)     & \BaseMaxTaskThin & \BaseMaxTaskMid & \BaseMaxTaskThick \\
Longest task, salted (s)       & \SaltMaxTaskThin & \SaltMaxTaskMid & \SaltMaxTaskThick \\
\textbf{Critical-path cut}     & \textbf{\SaltTaskCutThin$\times$} & \textbf{\SaltTaskCutMid$\times$} & \textbf{\SaltTaskCutThick$\times$} \\
\bottomrule
\end{tabular}
\end{center}

At $P=\PThin$ the result is inside the 18.3\% floor and we draw no conclusion.
At $P=\PMid$ and $P=\PThick$, where the floor is 0.1\% and 6.5\%, salting
returns \SaltSpeedupMid$\times$ and \SaltSpeedupThick$\times$ --- comfortably
outside noise --- and cuts the critical path by \SaltTaskCutMid$\times$ and
\SaltTaskCutThick$\times$ where AQE cut it by nothing.

One nuance worth recording, because it bears on the cost model. The $k$ that
minimises \emph{wall-clock} is \SaltBestKThick{} at $P=\PThick$, but a larger
$k$ minimises the \emph{critical path}: at $k=16$ the longest task falls
further still. The two optima differ because on two task slots the extra
shuffle from a larger $k$ costs more than the additional straggler relief
returns. On a cluster with more slots they should converge, and that is one of
the things the released cluster configuration is designed to check.

\subsection{Broadcast dominates at this scale, and we say so}

The broadcast arm completes in \BcastWallThin~s, \BcastWallMid~s and
\BcastWallThick~s --- faster than every shuffle-join strategy at every probe
thickness. This is expected and it is worth stating plainly rather than
omitting: our dimension table is well under Spark's default broadcast
threshold, and a practitioner in that regime should broadcast and stop reading.
Automatic broadcast is disabled throughout precisely so that the shuffle-join
behaviour we set out to study is not silently replaced. The results here apply
to joins whose dimension side exceeds the broadcast threshold, which is the
regime in which skew mitigation is a live question at all.

\subsection{Replication cost matches the model exactly}

Shuffle volume is a property of the plan and the data rather than of the
machine, so it is the measurement that transfers off a single node. Measured in
shuffle \emph{records} it does more than transfer: it matches
Equation~\eqref{eq:model} to the row.

An earlier version of this analysis measured replication in shuffle
\emph{megabytes} and reported a sevenfold gap between the two salting variants.
That was wrong, and the error is worth recording because it is easy to repeat.
Selective salting replicates only $P(k-1)$ rows --- a few hundred across the
entire sweep --- while the megabyte growth was dominated by the one-column
\texttt{\_salt} addition to every one of two million fact rows. Over 99\% of
what was attributed to replication was the extra column. Records have no such
confound: they are exactly the quantity the model's $bP(k-1)$ term refers to.

Fitting a least-squares line to shuffle records against $k$ over a matched
range gives:

\begin{center}
\begin{tabular}{lrrr}
\toprule
Variant & records per unit $k$ & model predicts & $R^{2}$ \\
\midrule
Selective salting & \SelRecSlope & $P = \PHOT$ & \SelRecR \\
Uniform salting   & \UnifRecSlope & $M = \NDIM$ & \UnifRecR \\
\bottomrule
\end{tabular}
\end{center}

Selective salting adds \SelRecSlope{} rows per salt value, and the hot key has
\PHOT{} rows on the probe side. Uniform salting adds \UnifRecSlope{} rows per
salt value, and the probe table has \NDIM{} rows. The measured ratio,
\RecSlopeRatio, equals $M/P$ exactly, which is what the model predicts, and
both fits are linear to $R^{2} = 1.000$.

This is the paper's cleanest quantitative result, and we would emphasise
\emph{why} it is clean rather than merely report it. Records are counted, not
timed. They do not vary between repetitions, they do not depend on the machine,
and they are not confounded by anything the scheduler does. The replication
term of the cost model is therefore not an approximation fitted to noisy
observations --- it is exact, and the distinction between the two salting
variants is a factor of $M/P$, which for any realistic fact--dimension join is
three orders of magnitude, not seven-fold.

The practical reading is that \textbf{the tutorial formulation of salting is
the expensive one}, and by a margin far wider than the folklore suggests.

\begin{table}[t]
\centering
\caption{Best configuration per strategy, with speedup over the unmitigated baseline. $\mu$ is the probe-side hot-key multiplier; skew is the join-stage maximum-to-median task duration ratio.}
\label{tab:headline}
\footnotesize
\begin{tabular}{lrrrrr}
\toprule
$\mu$ & strategy & $k$ & wall (s) & speedup & skew \\
\midrule
1.00 & baseline & 1 & 2.61 & 1.00 & 7.48 \\
1.00 & broadcast & 1 & 2.62 & 1.00 & 1.05 \\
1.00 & salt\_selective & 2 & 2.69 & 0.97 & 5.82 \\
1.00 & salt\_uniform & 2 & 2.78 & 0.94 & 4.97 \\
1.00 & aqe & 1 & 2.88 & 0.91 & 1.02 \\
1.00 & aqe\_salt & 8 & 3.01 & 0.87 & 2.50 \\
10.00 & salt\_selective & 2 & 5.19 & 1.34 & 23.87 \\
10.00 & aqe\_salt & 8 & 5.53 & 1.26 & 1.32 \\
10.00 & broadcast & 1 & 5.60 & 1.24 & 1.03 \\
10.00 & salt\_uniform & 8 & 5.68 & 1.23 & 6.39 \\
10.00 & aqe & 1 & 6.90 & 1.01 & 3.54 \\
10.00 & baseline & 1 & 6.96 & 1.00 & 35.97 \\
32.00 & salt\_selective & 2 & 11.08 & 1.68 & 68.38 \\
32.00 & aqe\_salt & 8 & 11.67 & 1.60 & 13.91 \\
32.00 & salt\_uniform & 8 & 11.72 & 1.59 & 17.45 \\
32.00 & broadcast & 1 & 13.32 & 1.40 & 1.03 \\
32.00 & aqe & 1 & 18.33 & 1.02 & 10.39 \\
32.00 & baseline & 1 & 18.62 & 1.00 & 101.01 \\
\bottomrule
\end{tabular}
\end{table}


\subsection{Where the cost model held, and where it did not}

The model's structural claims are supported. The replication term is exact, as
above. Salting monotonically reduces the critical path: at the thickest probe
side the longest join-stage task falls from \BaseWallThick~s unmitigated to a
small fraction of that under salting.

The \emph{convexity} of $T(k)$ in wall-clock, however, did not resolve, and we
report that plainly rather than presenting a fitted curve as a confirmation.

Fitting Equation~\eqref{eq:model} to the selective-salting sweeps gives
$R^{2}$ of \RsqThin, \RsqMid{} and \RsqThick. Those look respectable and are
misleading. The quantity that matters is whether the replication coefficient
$b$ --- the entire denominator of $k^{*}$ --- is resolved by the data at all:

\begin{center}
\begin{tabular}{lrrr}
\toprule
$P$ & \PThin & \PMid & \PThick \\
\midrule
$R^{2}$ & \RsqThin & \RsqMid & \RsqThick \\
$|b| / \mathrm{se}(b)$ & \BtstatThin & \BtstatMid & \BtstatThick \\
$b$ identified? & \IdentThin & \IdentMid & \IdentThick \\
$k^{*}$ & \KstarThin & \KstarMid & \KstarThick \\
\bottomrule
\end{tabular}
\end{center}

At two of three workloads $b$ is statistically indistinguishable from zero. A
$k^{*}$ computed from such a fit is an artefact of noise, and our
implementation refuses to return one: \texttt{optimal\_k\_unconstrained} yields
\textsc{NaN} rather than a number, and the recommendation falls back to the
saturation and parallelism ceilings alone. Only \NumIdentified{} of
\NumFits{} fits supports a cost-optimum at all.

A second, independent sign of trouble: the fitted $\gamma = a/b$ varies by
\GammaSpread$\times$ across three workloads that differ only in $P$, on one
machine, in one session. Section~\ref{sec:costmodel} advances $\gamma$-stability
as the property that would make the closed form practically useful. Our own
data contradict it.

We therefore make no claim that the closed form is validated. What we can say
is precise: the replication \emph{term} is exact; the replication \emph{cost in
time} is, at this scale, too small for the convexity it implies to be
observable; and the conditions under which it should become observable ---
cluster-scale network shuffle, or a substantially thicker probe side --- are
identified and configured in the released artifact.

The practical consequence points the opposite way from the folklore, and is
worth stating for practitioners rather than only for reviewers. \textbf{With
selective salting, over-provisioning $k$ is close to free.} The warning against
large $k$ is well founded for the uniform variant, where the replication cost
is $M/P$ times larger and clearly measurable. For the selective variant, the
sensible advice is to salt at the saturation bound and stop worrying about the
exact optimum.

\subsection{Summary}

\begin{itemize}
\item AQE's advantage over an unmitigated baseline fell from
\AqeSpeedupThin$\times$ to \AqeSpeedupThick$\times$ as the probe-side hot key
grew from \PThin{} to \PThick{} rows, while salting's rose from
\SaltSpeedupThin$\times$ to \SaltSpeedupThick$\times$.
\item Spark's skew-join rule never fired on a join with a
\BaseSkewThick:1 task-time skew ratio, verified from final adaptive plans. The
measured AQE benefit was partition coalescing.
\item Neither documented trigger knob was responsible; the binding condition
was \texttt{advisoryPartitionSizeInBytes} exceeding the hot partition. Lowering
it engaged the rule immediately.
\item Selective salting replicates \SelRecSlope{} rows per salt value against
uniform salting's \UnifRecSlope{} --- a ratio of \RecSlopeRatio, exactly the
$M/P$ the cost model predicts, with both fits linear to $R^{2}=1.000$.
\item The replication term of the cost model is confirmed exactly. Its
consequence for runtime is not: the coefficient $b$ is statistically
unidentified at two of three workloads, fitted $\gamma$ varies by
\GammaSpread$\times$ on one machine, and we therefore do not claim the closed
form for $k^{*}$ is validated.
\end{itemize}
